\documentclass[10pt,a4paper]{amsart}
\usepackage{amsmath,amssymb,amsthm,graphicx}
% Retained minimal preamble for the Stacks Project chapter flat.tex (More on Flatness).
% Verbatim excerpt of the chapter preamble: Xy-pic setup (native preamble lines 14-18) and the native
% theorem declarations with the mathematical macros (native preamble lines 154-217). Copied mechanically;
% nothing added, removed or redefined. The wrapper class and the amsmath/amssymb/amsthm/graphicx base
% packages are supplied by the assembler, so the class/externaldocument part of the native preamble is omitted.
% For commutative diagrams we use Xy-pic
\usepackage[all]{xy}

% We use 2cell for 2-commutative diagrams.
\xyoption{2cell}
% Theorem environments.
%
\theoremstyle{plain}
\newtheorem{theorem}[subsection]{Theorem}
\newtheorem{proposition}[subsection]{Proposition}
\newtheorem{lemma}[subsection]{Lemma}

\theoremstyle{definition}
\newtheorem{definition}[subsection]{Definition}
\newtheorem{example}[subsection]{Example}
\newtheorem{exercise}[subsection]{Exercise}
\newtheorem{situation}[subsection]{Situation}

\theoremstyle{remark}
\newtheorem{remark}[subsection]{Remark}
\newtheorem{remarks}[subsection]{Remarks}

\numberwithin{equation}{subsection}

% Macros
%
\def\lim{\mathop{\mathrm{lim}}\nolimits}
\def\colim{\mathop{\mathrm{colim}}\nolimits}
\def\Spec{\mathop{\mathrm{Spec}}}
\def\Hom{\mathop{\mathrm{Hom}}\nolimits}
\def\Ext{\mathop{\mathrm{Ext}}\nolimits}
\def\SheafHom{\mathop{\mathcal{H}\!\mathit{om}}\nolimits}
\def\SheafExt{\mathop{\mathcal{E}\!\mathit{xt}}\nolimits}
\def\Sch{\mathit{Sch}}
\def\Mor{\mathop{\mathrm{Mor}}\nolimits}
\def\Ob{\mathop{\mathrm{Ob}}\nolimits}
\def\Sh{\mathop{\mathit{Sh}}\nolimits}
\def\NL{\mathop{N\!L}\nolimits}
\def\CH{\mathop{\mathrm{CH}}\nolimits}
\def\proetale{{pro\text{-}\acute{e}tale}}
\def\etale{{\acute{e}tale}}
\def\QCoh{\mathit{QCoh}}
\def\Ker{\mathop{\mathrm{Ker}}}
\def\Im{\mathop{\mathrm{Im}}}
\def\Coker{\mathop{\mathrm{Coker}}}
\def\Coim{\mathop{\mathrm{Coim}}}

% Boxtimes
%
\DeclareMathSymbol{\boxtimes}{\mathbin}{AMSa}{"02}

%
% Macros for moduli stacks/spaces
%
\def\QCohstack{\mathcal{QC}\!\mathit{oh}}
\def\Cohstack{\mathcal{C}\!\mathit{oh}}
\def\Spacesstack{\mathcal{S}\!\mathit{paces}}
\def\Quotfunctor{\mathrm{Quot}}
\def\Hilbfunctor{\mathrm{Hilb}}
\def\Curvesstack{\mathcal{C}\!\mathit{urves}}
\def\Polarizedstack{\mathcal{P}\!\mathit{olarized}}
\def\Complexesstack{\mathcal{C}\!\mathit{omplexes}}
% \Pic is the operator that assigns to X its picard group, usage \Pic(X)
% \Picardstack_{X/B} denotes the Picard stack of X over B
% \Picardfunctor_{X/B} denotes the Picard functor of X over B
\def\Pic{\mathop{\mathrm{Pic}}\nolimits}
\def\Picardstack{\mathcal{P}\!\mathit{ic}}
\def\Picardfunctor{\mathrm{Pic}}
\def\Deformationcategory{\mathcal{D}\!\mathit{ef}}
\usepackage{cleveref}
\crefname{originaldefinition}{original definition}{original definitions}
\Crefname{originaldefinition}{Original definition}{Original definitions}
\makeatletter
\newlabel{situation-existence-derived}{{0DIB}{source}{}{sourcepointer-1}{}}
\newlabel{situation-existence-derived@cref}{{[originaldefinition][0][]0DIB}{[1][0][]source}{}{}{}}
\makeatother
\crefname{originallemma}{original lemma}{original lemmas}
\Crefname{originallemma}{Original lemma}{Original lemmas}
\makeatletter
\newlabel{perfect-lemma-better-coherator}{{0CR0}{source}{}{sourcepointer-2}{}}
\newlabel{perfect-lemma-better-coherator@cref}{{[originallemma][0][]0CR0}{[1][0][]source}{}{}{}}
\makeatother
\crefname{originallemma}{original lemma}{original lemmas}
\Crefname{originallemma}{Original lemma}{Original lemmas}
\makeatletter
\newlabel{perfect-lemma-boundedness-better-coherator}{{0CSA}{source}{}{sourcepointer-3}{}}
\newlabel{perfect-lemma-boundedness-better-coherator@cref}{{[originallemma][0][]0CSA}{[1][0][]source}{}{}{}}
\makeatother
\crefname{originallemma}{original lemma}{original lemmas}
\Crefname{originallemma}{Original lemma}{Original lemmas}
\makeatletter
\newlabel{cohomology-lemma-RGamma-commutes-with-Rlim}{{0D60}{source}{}{sourcepointer-4}{}}
\newlabel{cohomology-lemma-RGamma-commutes-with-Rlim@cref}{{[originallemma][0][]0D60}{[1][0][]source}{}{}{}}
\makeatother
\crefname{originallemma}{original lemma}{original lemmas}
\Crefname{originallemma}{Original lemma}{Original lemmas}
\makeatletter
\newlabel{perfect-lemma-pseudo-coherent}{{08E5}{source}{}{sourcepointer-5}{}}
\newlabel{perfect-lemma-pseudo-coherent@cref}{{[originallemma][0][]08E5}{[1][0][]source}{}{}{}}
\makeatother
\crefname{originallemma}{original lemma}{original lemmas}
\Crefname{originallemma}{Original lemma}{Original lemmas}
\makeatletter
\newlabel{perfect-lemma-affine-compare-bounded}{{06Z0}{source}{}{sourcepointer-6}{}}
\newlabel{perfect-lemma-affine-compare-bounded@cref}{{[originallemma][0][]06Z0}{[1][0][]source}{}{}{}}
\makeatother
\crefname{originallemma}{original lemma}{original lemmas}
\Crefname{originallemma}{Original lemma}{Original lemmas}
\makeatletter
\newlabel{cohomology-lemma-finite-cohomology}{{08DN}{source}{}{sourcepointer-7}{}}
\newlabel{cohomology-lemma-finite-cohomology@cref}{{[originallemma][0][]08DN}{[1][0][]source}{}{}{}}
\makeatother
\crefname{originallemma}{original lemma}{original lemmas}
\Crefname{originallemma}{Original lemma}{Original lemmas}
\makeatletter
\newlabel{algebra-lemma-NAK}{{00DV}{source}{}{sourcepointer-8}{}}
\newlabel{algebra-lemma-NAK@cref}{{[originallemma][0][]00DV}{[1][0][]source}{}{}{}}
\makeatother
\title{The derived limit of the quasi-coherent approximations lies in bounded above quasi-coherent complexes}
\author{The Stacks Project Authors}
\date{stacks-project:flat.tex:local-snapshot-20261011; GNU FDL 1.2 or later}
\begin{document}
\maketitle
\noindent\textit{Editorial note.} All delivered mathematical text is the author's own, copied from the pinned Stacks Project chapter file flat.tex; no statement or proof step was rewritten, reordered or removed.\par\medskip
\noindent\textit{Editorial note.} The preamble is the chapter's own retained minimal preamble (sources/stacks-flat/retained-preamble.tex), a verbatim excerpt of the native Xy-pic setup and of the native theorem declarations and mathematical macros; no macro was invented or redefined, and the wrapper class and the amsmath/amssymb/amsthm/graphicx base packages are supplied by the assembler.\par\medskip
\noindent\textit{Editorial note.} The counted statement and proof contain no cross reference, no citation and no included file, so the delivered body is self contained.\par\medskip
\noindent\textit{Editorial note.} Canonical location of the counted proposition: Stacks tag 0DIC (sources/stacks-flat/pinned-original-source.tex line 8406).\par\medskip
\noindent\textit{Allowed standard prerequisites (editorial).}\begin{itemize}
\item Existence of the coherator as a right adjoint with bounded cohomological dimension. Derived Categories of Schemes, Lemmas labelled perfect-lemma-better-coherator (Stacks tag 0CR0) and perfect-lemma-boundedness-better-coherator (Stacks tag 0CSA), invoked at proof lines 8417 and 8424-8426. Supply the existence of DQ\_X and its commutation with derived limits, and the bounded cohomological dimension that reduces the claim to R-lim K\_n in D\textasciicircum{}-(O\_X).
\item The exact sequence relating the derived limit of cohomology to the cohomology of the derived limit, and the comparison of cohomology for pseudo-coherent complexes on an affine. Cohomology, Lemma labelled cohomology-lemma-RGamma-commutes-with-Rlim (Stacks tag 0D60) at proof lines 8429-8435, and Derived Categories of Schemes, Lemmas labelled perfect-lemma-pseudo-coherent (Stacks tag 08E5) and perfect-lemma-affine-compare-bounded (Stacks tag 06Z0) at proof lines 8436-8441. Supply the reduction of the claim to the boundedness of the K\_n independent of n.
\item The situation of the existence-of-hulls setup and Nakayama for a finite presentation module along a thickening. Same chapter, Situation labelled situation-existence-derived (source line 8379) invoked in the statement, Cohomology, Lemma labelled cohomology-lemma-finite-cohomology (Stacks tag 08DN) at proof line 8451 and Algebra, Lemma labelled algebra-lemma-NAK (Stacks tag 00DV) at proof lines 8454-8456. Supply the tower X\_n, the pseudo-coherence of K\_n and the finite presentability that makes Nakayama applicable; the comparison of the degrees a\_n is carried out in the counted span.
\end{itemize}\medskip
\ifdefined\hypertarget\hypertarget{sourcepointer-1}{}\fi
\noindent\textit{Source reference.} Original definition 0DIB: sources/stacks-flat/pinned-original-source.tex, source lines 8379--8379. Same-chapter pointer into the author environment situation at source line 8379; invoked at proof line None of the counted proof. It is retained only as a navigation pointer; every load-bearing step is inside the delivered span.\par\medskip
\ifdefined\hypertarget\hypertarget{sourcepointer-2}{}\fi
\noindent\textit{Source reference.} Original lemma 0CR0: sources/stacks-chow/tags.txt, source lines 15980--15980. Cross-chapter standard prerequisite located in the project tag list staged in the local snapshot; invoked at proof line 8418 of the counted proof. It is retained only as a navigation pointer; every load-bearing step is inside the delivered span.\par\medskip
\ifdefined\hypertarget\hypertarget{sourcepointer-3}{}\fi
\noindent\textit{Source reference.} Original lemma 0CSA: sources/stacks-chow/tags.txt, source lines 16025--16025. Cross-chapter standard prerequisite located in the project tag list staged in the local snapshot; invoked at proof line 8425 of the counted proof. It is retained only as a navigation pointer; every load-bearing step is inside the delivered span.\par\medskip
\ifdefined\hypertarget\hypertarget{sourcepointer-4}{}\fi
\noindent\textit{Source reference.} Original lemma 0D60: sources/stacks-chow/tags.txt, source lines 16505--16505. Cross-chapter standard prerequisite located in the project tag list staged in the local snapshot; invoked at proof line 8435 of the counted proof. It is retained only as a navigation pointer; every load-bearing step is inside the delivered span.\par\medskip
\ifdefined\hypertarget\hypertarget{sourcepointer-5}{}\fi
\noindent\textit{Source reference.} Original lemma 08E5: sources/stacks-chow/tags.txt, source lines 10661--10661. Cross-chapter standard prerequisite located in the project tag list staged in the local snapshot; invoked at proof line 8438 of the counted proof. It is retained only as a navigation pointer; every load-bearing step is inside the delivered span.\par\medskip
\ifdefined\hypertarget\hypertarget{sourcepointer-6}{}\fi
\noindent\textit{Source reference.} Original lemma 06Z0: sources/stacks-chow/tags.txt, source lines 8904--8904. Cross-chapter standard prerequisite located in the project tag list staged in the local snapshot; invoked at proof line 8441 of the counted proof. It is retained only as a navigation pointer; every load-bearing step is inside the delivered span.\par\medskip
\ifdefined\hypertarget\hypertarget{sourcepointer-7}{}\fi
\noindent\textit{Source reference.} Original lemma 08DN: sources/stacks-chow/tags.txt, source lines 10643--10643. Cross-chapter standard prerequisite located in the project tag list staged in the local snapshot; invoked at proof line 8451 of the counted proof. It is retained only as a navigation pointer; every load-bearing step is inside the delivered span.\par\medskip
\ifdefined\hypertarget\hypertarget{sourcepointer-8}{}\fi
\noindent\textit{Source reference.} Original lemma 00DV: sources/stacks-chow/tags.txt, source lines 512--512. Cross-chapter standard prerequisite located in the project tag list staged in the local snapshot; invoked at proof line 8455 of the counted proof. It is retained only as a navigation pointer; every load-bearing step is inside the delivered span.\par\medskip
\begin{lemma}
\label{lemma-compute-what-it-should-be-derived}
In Situation \ref{situation-existence-derived} consider
$$
K = R\lim_{D_\QCoh(\mathcal{O}_X)}(K_n) =
DQ_X(R\lim_{D(\mathcal{O}_X)} K_n)
$$
Then $K$ is in $D^-_{\QCoh}(\mathcal{O}_X)$.
\end{lemma}

\begin{proof}
The functor $DQ_X$ exists because $X$ is quasi-compact and
quasi-separated, see Derived Categories of Schemes, Lemma
\ref{perfect-lemma-better-coherator}.
Since $DQ_X$ is a right adjoint it commutes with products
and therefore with derived limits. Hence the equality
in the statement of the lemma.

\medskip\noindent
By Derived Categories of Schemes,
Lemma \ref{perfect-lemma-boundedness-better-coherator}
the functor $DQ_X$ has bounded cohomological dimension.
Hence it suffices to show that $R\lim K_n \in D^-(\mathcal{O}_X)$.
To see this, let $U \subset X$ be an affine open.
Then there is a canonical exact sequence
$$
0 \to
R^1\lim H^{m - 1}(U, K_n) \to H^m(U, R\lim K_n) \to
\lim H^m(U, K_n) \to 0
$$
by Cohomology, Lemma \ref{cohomology-lemma-RGamma-commutes-with-Rlim}.
Since $U$ is affine and $K_n$ is pseudo-coherent (and hence has
quasi-coherent cohomology sheaves by
Derived Categories of Schemes, Lemma \ref{perfect-lemma-pseudo-coherent})
we see that $H^m(U, K_n) = H^m(K_n)(U)$ by
Derived Categories of Schemes, Lemma
\ref{perfect-lemma-affine-compare-bounded}.
Thus we conclude that it suffices to show that $K_n$
is bounded above independent of $n$.

\medskip\noindent
Since $K_n$ is pseudo-coherent we have $K_n \in D^-(\mathcal{O}_{X_n})$.
Suppose that $a_n$ is maximal such that $H^{a_n}(K_n)$ is nonzero.
Of course $a_1 \leq a_2 \leq a_3 \leq \ldots$.
Note that $H^{a_n}(K_n)$ is an
$\mathcal{O}_{X_n}$-module of finite presentation
(Cohomology, Lemma \ref{cohomology-lemma-finite-cohomology}).
We have $H^{a_n}(K_{n - 1}) =
H^{a_n}(K_n) \otimes_{\mathcal{O}_{X_n}} \mathcal{O}_{X_{n - 1}}$.
Since $X_{n - 1} \to X_n$ is a thickening, it follows from
Nakayama's lemma (Algebra, Lemma \ref{algebra-lemma-NAK}) that if
$H^{a_n}(K_n) \otimes_{\mathcal{O}_{X_n}} \mathcal{O}_{X_{n - 1}}$
is zero, then $H^{a_n}(K_n)$ is zero too. Thus $a_n = a_{n - 1}$
for all $n$ and we conclude.
\end{proof}

\end{document}
